\scanpage{12}
\deleted{Localisé d'un schéma en un point, schéma local, morphisme d'un schéma local dans un schéma.}

\deleted{I Morphismes quasi-compacts de type fini loc. de type fini loc. de prés. finie}

\diagram{\begin{tikzcd}[ampersand replacement=\&]
 \deleted{X}\arrow[d,"f"']\&\deleted{X'}\arrow[l]\arrow[d,"f'"]\\
 \deleted{Y}\arrow[r,dash]\&\deleted{Y'}
\end{tikzcd}}

\diagram{\begin{tikzcd}[ampersand replacement=\&]
 \deleted{X}\arrow[rr,"f"]\arrow[dr,dash]\&{}\&\deleted{Y}\arrow[dl,dash,"g"]\\
 {}\&\deleted{S}\&{}
\end{tikzcd}}

\deleted{Si $f$ est $M$, $f'$ est $M$. Si $f$ et $g$ sont $M$, $fg$ aussi $\Rightarrow$ Si $f_1$ et $f_2$ sont $M$, $f_1\times_S f_2$ est $M$}

\deleted{Si $gf$ est $M$, $f$ est $M$] vrai pour \unsure{loc. de t.f.|loc. t.f.}, \deleted{\unsure{loc. prés.}} imm., morphismes séparés}

\deleted{immersions $\left\{\begin{aligned}&\text{ouvertes}\\&\text{fermées}\end{aligned}\right.$}

\margin{\deleted{immersion surjective \illegible{4} \unsure{ou} \illegible{2}}}

\deleted{Sous-schémas ouverts -- fermés}

\deleted{\uline{Ex} \{ Morphisme diagonal -- noyaux \quad sections}

\deleted{morphismes séparés : diagonale fermée -- noyaux fermés -- principe de prolongement des identités $Z\longrightarrow X$ pour $Z$ réduit ($X$ séparé sur $S$)}

\deleted{morphismes qu}

\textcircled{II} $\left\{\begin{aligned}
 &\text{Schémas noeth, loc. noeth, artiniens, loc. artiniens}\\
 &\text{Schémas irréd, connexes}\\
 &\text{Schémas \deleted{réduits}, intègres}\\
 &\text{Schémas normaux, réguliers}
\end{aligned}\right.$

I) \unsure{Schéma localisé|Schémas localisé|Schémas localisés} d'un \deleted{\unsure{sc}} schéma en un point schémas locaux ($\exists!$ pt fermé, et quasi-compact) Morphisme d'un schéma \inserted{local} dans un schéma.

\textcircled{III} Morphismes : \uline{affines}, \inserted{[mono, surjectifs]} [\uline{qu. cpt}, \uline{de t.f}, \uline{loc.} de t.f., \uline{loc.} de \uline{p.f.}], \uline{imm.} (\uline{ouv.}, \uline{fermées}, \uline{surjectives}), \uline{séparés}, [\uline{quasi-séparés} -- \uline{cohérents}, de \uline{prés. finie}]

\textcircled{c} \uline{morphismes} \inserted{[\uline{entiers}]} \uline{finis}, [quasi-finis]

\deleted{\unsure{quasi}} \textcircled{a} (\uline{projectifs}, \uline{quasi-projectifs}) [quasi-affines]

\textcircled{d} \uline{radiciels} homéom. universels \deleted{\unsure{univ|univ.}} ouverts, univ. fermés

\textcircled{b} \uline{propres}

\margin{Sur base cohérente}

\uline{Oubli} : \{ Cohérence du faisceau structural sur un schéma loc. noeth, \uline{structure induite réduite} sur un \uline{fermé}

\inserted{attention, une imm. ouverte n'est pas une imm. \unsure{quasi entière|quasi-entière}, \unsure{idée} \unsure{mettre} : cas des deux \unsure{variétés même corps|valuations même corps})}

\inserted{où on rencontre les schémas non séparés}

Un schéma loc. noeth. est qu. sép.

$f$ cohérent, $F$ qu. coh $\Rightarrow f_*(F)$ qu. coh.

\uline{cor} \{ Extension de faisceaux qu. coh sur ouvert rétrocompact

Extension de sous-faisceaux \deleted{\unsure{de t.f. de type}} \unsure{de} $F|U$ ($U$ rétro-cpt)

$\left\{\begin{aligned}
 &\text{\uline{net}, \uline{lisse}, \uline{étale}}\\
 &\boxed{\text{\uline{plat}, \uline{fid plat}}}
\end{aligned}\right.$

Propriétés géom des fibres \ldots

% archive sheet number: 11
